\section*{Розділ VII. Взаємодія з іншими органами}
\addcontentsline{toc}{section}{Розділ VII. Взаємодія з іншими органами}

\subsection*{7.1. Взаємодія з Конференцією студентів Університету (КСУ)}
\addcontentsline{toc}{subsection}{7.1. Взаємодія з Конференцією студентів Університету (КСУ)}
    7.1.1. ЦВКс підзвітна КСУ та надає їй звіти про свою діяльність у встановленому порядку.

    7.1.2. ЦВКс взаємодіє з КСУ з питань:

        \begin{enumerate}[label=\alph*)]
            \item Затвердження загальних принципів виборчих процедур;
            \item Передачі результатів виборів Голови СР КАІ та Голови СР СМ для затвердження;
            \item Організаційно-технічної підтримки виборів до СУД;
            \item Розгляду скарг на рішення ЦВКс;
            \item Внесення змін до виборчих процедур.
        \end{enumerate}

    7.1.3. КСУ може давати доручення ЦВКс щодо проведення конкретних виборів та встановлення додаткових вимог до виборчих процедур.

\subsection*{7.2. Взаємодія зі Студентською радою КАІ (СР КАІ)}
\addcontentsline{toc}{subsection}{7.2. Взаємодія зі Студентською радою КАІ (СР КАІ)}
    7.2.1. СР КАІ сприяє ЦВКс в організації та проведенні виборів до органів студентського самоврядування Університету.

    7.2.2. Співпраця ЦВКс та СР КАІ включає:

        \begin{enumerate}[label=\alph*)]
            \item Надання ЦВКс необхідної інформації для організації виборчого процесу;
            \item Сприяння в інформаційному супроводі виборчих кампаній через інформаційні ресурси СР КАІ;
            \item Надання приміщень або технічних ресурсів для потреб ЦВКс під час виборчого процесу;
            \item Координацію роботи з СРФ/СРІ щодо організації виборів на факультетському рівні.
        \end{enumerate}

    7.2.3. СР КАІ не втручається у здійснення ЦВКс її виключних повноважень щодо організації та проведення виборів.

\subsection*{7.3. Взаємодія зі Студентською уповноваженою делегацією (СУД)}
\addcontentsline{toc}{subsection}{7.3. Взаємодія зі Студентською уповноваженою делегацією (СУД)}
    7.3.1. ЦВКс взаємодіє з СУД з питань забезпечення законності виборчого процесу та розгляду порушень виборчого законодавства.

    7.3.2. ЦВКс має право звертатися до СУД з поданнями про:

        \begin{enumerate}[label=\alph*)]
            \item Порушення виборчого законодавства учасниками виборчого процесу;
            \item Необхідність застосування дисциплінарних санкцій до порушників;
            \item Надання роз'яснень щодо застосування норм виборчого права.
        \end{enumerate}

    7.3.3. СУД може надавати ЦВКс консультації з правових питань виборчого процесу та брати участь у розробці виборчих процедур як експертний орган.

\subsection*{7.4. Взаємодія зі студентськими радами факультетів/інститутів (СРФ/СРІ)}
\addcontentsline{toc}{subsection}{7.4. Взаємодія зі студентськими радами факультетів/інститутів (СРФ/СРІ)}
    7.4.1. ЦВКс координує з СРФ/СРІ проведення виборів на факультетському/інститутському рівні.

    7.4.2. СРФ/СРІ зобов'язані:

        \begin{enumerate}[label=\alph*)]
            \item Надавати ЦВКс інформацію про кількість студентів та структуру факультету/інституту;
            \item Сприяти в організації виборчих дільниць на території факультету/інституту;
            \item Забезпечувати інформування студентів про виборчий процес;
            \item Виконувати рішення ЦВКс з питань організації виборів.
        \end{enumerate}

    7.4.3. ЦВКс може делегувати СРФ/СРІ окремі технічні функції з організації виборів на факультетському рівні під своїм контролем.

\subsection*{7.5. Взаємодія з адміністрацією Університету}
\addcontentsline{toc}{subsection}{7.5. Взаємодія з адміністрацією Університету}
    7.5.1. Взаємодія ЦВКс з адміністрацією Університету здійснюється на принципах взаємоповаги, незалежності ЦВКс у виборчих питаннях та сприяння демократичному розвитку студентського самоврядування.

    7.5.2. Адміністрація Університету сприяє діяльності ЦВКс шляхом:

        \begin{enumerate}[label=\alph*)]
            \item Надання приміщень для проведення виборів та засідань ЦВКс;
            \item Забезпечення технічної підтримки виборчого процесу;
            \item Надання інформації про кількість студентів та їх розподіл по підрозділах;
            \item Сприяння в оприлюдненні інформації про вибори через офіційні канали Університету.
        \end{enumerate}

    7.5.3. Адміністрація Університету не втручається у виборчий процес та рішення ЦВКс, крім випадків порушення законодавства України.

\subsection*{7.6. Взаємодія зі студентськими організаціями}
\addcontentsline{toc}{subsection}{7.6. Взаємодія зі студентськими організаціями}
    7.6.1. ЦВКс взаємодіє зі студентськими організаціями як з суб'єктами виборчого процесу, які мають право висувати кандидатів та брати участь у виборчій кампанії.

    7.6.2. Студентські організації можуть:

        \begin{enumerate}[label=\alph*)]
            \item Висувати кандидатів на виборні посади в органах студентського самоврядування;
            \item Направляти своїх представників як спостерігачів на вибори;
            \item Брати участь у передвиборчій агітації відповідно до встановлених правил;
            \item Звертатися до ЦВКс зі скаргами та зауваженнями щодо виборчого процесу.
        \end{enumerate}

    7.6.3. ЦВКс забезпечує рівні права та можливості для всіх студентських організацій у виборчому процесі незалежно від їх спрямування та розміру.

\subsection*{7.7. Міжнародна співпраця}
\addcontentsline{toc}{subsection}{7.7. Міжнародна співпраця}
    7.7.1. ЦВКс може розвивати співпрацю з органами студентського самоврядування інших вищих навчальних закладів України та зарубіжних країн з метою обміну досвідом організації виборчих процесів.

    7.7.2. ЦВКс може брати участь у міжнародних програмах, конференціях та семінарах з питань розвитку студентської демократії та виборчих технологій.

    7.7.3. Міжнародна діяльність ЦВКс здійснюється у координації з адміністрацією Університету та відповідними органами студентського самоврядування.